\begin{table}[ht]
\centering
\scriptsize
\renewcommand{\arraystretch}{1.5}
\caption{Benchmark comparison of numerical methods: Shooting, SciPy collocation-style BVP, and finite difference}
\label{tab:method_comparison_table}
\begin{adjustbox}{max width=\textwidth}
\begin{tabular}{|>{\centering\arraybackslash}m{2.8cm}|p{5.2cm}|p{5.2cm}|}
    \hline
    \textbf{Method} & \textbf{Pros} & \textbf{Cons} \\
    \hline
    \textbf{Shooting Method} &
    - Simple baseline implementation \newline
    - Useful for diagnosing IVP sensitivity \newline
    - Reuses standard IVP solvers
    &
    - Struggles with stiff ODE systems \newline
    - Sensitive to initial guesses \newline
    - May fail near thermodynamic pinch \\
    \hline
    \textbf{SciPy BVP} &
    - Adaptive mesh and residual control \newline
    - Current reference path in the benchmark \newline
    - Directly records solver metadata
    &
    - Depends on initialization and scaling \newline
    - Can be slower than shooting on easy cases \newline
    - Convergence is case dependent \\
    \hline
    \textbf{Finite Difference Method} &
    - Transparent spatial discretization \newline
    - Simple to compare against other methods \newline
    - Works with mixed boundary conditions
    &
    - Computationally intensive \newline
    - Requires solving large nonlinear systems \newline
    - More difficult to scale for large grids \\
    \hline
\end{tabular}
\end{adjustbox}
\end{table}
